\documentclass[openany,12pt,a4paper]{report}
\usepackage{subfiles}
\usepackage[]{graphicx}
\graphicspath{{./img/}{./../img/}}
\usepackage{../StileDoc}
\title{Studio di fattibilità}
\author{}
\date{\today}

\loadglsentries{../Glossario/Definizioni}

%Ultima versione documento

\newcommand{\versione}{1.0.0}

\begin{document}
	\makeatletter
	\begin{titlepage}
		\setlength{\headsep}{0pt}
		\begin{center}
			\includegraphics[width=0.5\linewidth]{Logo.png}\\[1em]
			{\huge \bfseries  \@title }\\[10ex]
			\textbf{\Large Informazioni Documento} \\[2em]
			\bgroup
			\def\arraystretch{1.5}
			\begin{tabular}{l|l}
				\textbf{Versione} & \versione{} \\
				\textbf{Data approvazione} & \large \@date \\
				\textbf{Responsabile} & Marco Focchiatti \\
				\textbf{Redattori} & Giulio Rossetti \\
				\textbf{Verificatori} & Matteo Rizzo \\
				\textbf{Distribuzione} & Prof. Tullio Vardanega \\
				& Prof. Riccardo Cardin \\
				& Gruppo Graphite \\
				\textbf{Uso} & Interno \\
				\textbf{Recapito} & graphite.swe@gmail.com \\
			\end{tabular}
			\egroup
		\end{center}
	\end{titlepage}
	\makeatother

	\thispagestyle{empty}
	\newpage
	
	%REGISTRO DELLE MODIFICHE
	
	\chapter*{Registro delle modifiche}
	\begin{center}	
		\begin{tabular}{|c|c|c|c|p{4cm}|}
			\hline
			\textbf{Versione} & \textbf{Data} & \textbf{Autore} & \textbf{Ruolo} & \textbf{Modifiche} \\ \hline
			
		\end{tabular}
	\end{center}

	% INDICE

	\tableofcontents

	% INTRODUZIONE

	\chapter{Introduzione}

	\section{Scopo}

	Lo studio di fattibilità ha l’obiettivo di descrivere le motivazioni che hanno portato
	il gruppo alla scelta del \glossario{capitolato}{capitolato} C3 e alla relativa esclusione degli altri capitolati proposti.

	\section{Glossario}

	Al fine di evitare ogni ambiguità relativa al linguaggio utilizzato nei documenti, viene fornito il \textit{Glossario v1.0.0}, contenente la definizione dei termini in corsivo marcati con la lettera "G" in pedice.

	\section{Riferimenti}

	\subsection{Riferimenti normativi}

	\begin{itemize}

		\item \textbf{Norme di progetto:} documento \textit{Norme di progetto v1.0.0}.

	\end{itemize}

	\subsection{Riferimenti informativi}

	\begin{itemize}
		\item \textbf{Capitolato d'appalto C1:} AJarvis - Assistente Virtuale Cerimonie Agile \\ \url{http://www.math.unipd.it/~tullio/IS-1/2017/Progetto/C1.pdf}

		\item \textbf{Capitolato d'appalto C2:} BlockCV: blockchain per gestione di CV certificati \\ \url{http://www.math.unipd.it/~tullio/IS-1/2017/Progetto/C2.pdf}

		\item \textbf{Capitolato d'appalto C3:} DeSpeect: interfaccia grafica per Speect \\ \url{http://www.math.unipd.it/~tullio/IS-1/2017/Progetto/C3.pdf}

		\item \textbf{Capitolato d'appalto C4:} ECoRe: Enterprise Content Recommendation \\ \url{http://www.math.unipd.it/~tullio/IS-1/2017/Progetto/C4.pdf}

		\item \textbf{Capitolato d'appalto C5:} Ironworks \\ \url{http://www.math.unipd.it/~tullio/IS-1/2017/Progetto/C5.pdf}

		\item \textbf{Capitolato d'appalto C6:} Marvin: dimostratore di Uniweb su Ethereum \\ \url{http://www.math.unipd.it/~tullio/IS-1/2017/Progetto/C6.pdf}

		\item \textbf{Capitolato d'appalto C7:} OpenAPM: cruscotto di Application Performance Management \\ \url{http://www.math.unipd.it/~tullio/IS-1/2017/Progetto/C7.pdf}

		\item \textbf{Capitolato d'appalto C8:} TuTourSelf: piattaforma di prenotazioni per artisti in tournee \\ \url{http://www.math.unipd.it/~tullio/IS-1/2017/Progetto/C8.pdf}
	\end{itemize}

	% CAPITOLATO SCELTO

	\chapter{Capitolato scelto}

	\section{Capitolato C3}

	\begin{itemize}
		\item \textbf{Nome:} \textit{DeSpeect}: interfaccia grafica per Speect;
		\item \textbf{Proponente:} MIVOQ S.R.L.;
		\item \textbf{Commitente:} Prof.Tullio Vardanega e Prof. Riccardo Cardin.
	\end{itemize}

	\subsection{Descrizione}

	Realizzazione di un'interfaccia grafica per \glossario{Speect}{Speect} [Meraka Institute(2008-2013)], una libreria per la creazione di sistemi di sintesi vocale, che agevoli l’ispezione del suo stato interno durante il funzionamento e la scrittura di test per le sue funzionalità.


	\subsection{Dominio applicativo}

	Il progetto si prefigge di fornire uno strumento di supporto all'utilizzo della libreria \textit{Speect}.

	\subsection{Dominio tecnologico}

	\begin{itemize}
		\item \textbf{interfaccia grafica:} \textit{suite Qt, Speect, C++, json, gnu GCC.}
	\end{itemize}

	\subsection{Aspetti positivi}

	\begin{itemize}
		\item la libreria ha uno scopo molto rilevante;

		\item il prodotto da realizzare può semplificare molto l'utilizzo di \textit{Speect};

		\item alcune tecnologie richieste dal progetto sono già familiari ai membri del gruppo e apparentemente piuttosto diffuse nell'ambito di pertinenza;

		\item il prodotto da realizzare ha prospettive di grande adattabilità all'accrescersi e all'espandersi della libreria cui fornisce supporto.
	\end{itemize}

	\subsection{Fattori di rischio}

	\begin{itemize}
		\item A seguito di un'analisi iniziale, la libreria sembra non essere particolarmente ben documentata;

		\item L'istallazione della libreria è ostica e complessa.
	\end{itemize}

	\subsection{Valutazione finale}

	Dall'esito dell'analisi del rapporto tra aspetti positivi e rischi concernenti la proposta di progetto e dal grande interesse che ha suscitato nei menbri del gruppo, la nostra scelta ricade sul capitolato C3.

	\chapter{Analisi degli altri capitolati}

	% CAPITOLATO 1

	\section{Capitolato C1}

	\begin{itemize}
		\item \textbf{Nome:} \textit{AJarvis} - Assistente Virtuale Cerimonie Agile;
		\item \textbf{Proponente:} Zero12;
		\item \textbf{Commitente:} Prof.Tullio Vardanega e Prof. Riccardo Cardin.
	\end{itemize}

	\subsection{Descrizione}

	Produzione di un’applicazione di \textit{Machine Learning} in grado di ascoltare gli stand up giornalieri sullo stato di avanzamento dei progetti, comprenderne i dialoghi e analizzarne il contenuto per fornire una sintesi dello stand up estraendo dal contesto gli argomenti emersi.

	\subsection{Dominio applicativo}

	Il progetto si inserisce nell'ambito dell'automatizzazione della creazione di report sulle riunioni stand up aziendali nel contesto della progettazione \glossario{Agile}{Agile}.

	\subsection{Dominio tecnologico}

	\begin{itemize}
		\item \textbf{Interfacce Web di registrazione e di reportistica delle analisi:} \textit{HTML5, CSS3, Javascript, AngularJS, PHP};

		\item \textbf{Servizi Cloud:} tecnologie \textit{Google Cloud Platform} per l’analisi dei dati e traduzione di voce in testo.
	\end{itemize}

	\subsection{Aspetti positivi}

	\begin{itemize}
		\item Proposta d'uso di tecnologie moderne e innovative;

		\item Possibilità di applicazione in un ambito aziendale concreto.
	\end{itemize}

	\subsection{Fattori di rischio}

	\begin{itemize}
		\item Complessità elevata nel produrre un prodotto di qualità utilizzando le tecnologie richieste, in particolare nel caso vengano riconosciute più voci simultaneamente, nonostante venga considerata come pre-condizione che ogni partecipante parli uno alla volta, può capitare durante una conversazione che inavvertitamente più persone parlino contemporaneamente;

		\item Il gruppo ritiene significativo che un prodotto così interessante sia condiviso sotto licenze \glossario{open source}{open source}, contrariamente a quanto sottolineato dalla proponente.
	\end{itemize}

	\subsection{Valutazione finale}

	Nonostante l'interesse suscitato dalle tecnologie proposte, la gravosità dei fattori di rischio hanno spinto il gruppo alla scelta di un altro capitolato.

	% CAPITOLATO 2

	\section{Capitolato C2}

	\begin{itemize}
		\item \textbf{Nome:} \textit{BlockCV}: blockchain per gestione di CV certificati;
		\item \textbf{Proponente:} Ifin Sistemi Srl;
		\item \textbf{Commitente:} Prof.Tullio Vardanega e Prof. Riccardo Cardin.
	\end{itemize}

	\subsection{Descrizione}

	Creazione di un sistema distribuito per la pubblicazione di CV certificati e per la ricerca di proposte di lavoro basato su una \textit{permissioned blockchain}.

	\subsection{Dominio applicativo}

	Si vuole ottenere un modello applicabile nell'attuale sistema lavorativo in grado di gestire le operazioni di base che accompagnano il lavoratore dalla fase iniziale di ricerca di una prima occupazione, con conseguente creazione e pubblicazione di un CV, seguendone via via l’evoluzione della carriera, attraverso i vari aggiornamenti del CV certificati dalle aziende in cui il lavoratore ha prestato servizio.

	\subsection{Dominio tecnologico}

	\begin{itemize}
		\item \textbf{Piattaforma blockchain:} \textit{Hyperledger Fabric, MongoDB, Cassandra};

		\item \textbf{Codifica:} \textit{Java EE};

		\item \textbf{Interfaccia grafica:} \textit{Framework Play} o suite di componenti \textit{Vaadin Elements}.
	\end{itemize}

	\subsection{Aspetti positivi}

	\begin{itemize}
		\item Realizzazione del progetto basata sull'utilizzo di una tecnologia innovativa e in forte sviluppo;

		\item Obiettivo del progetto interessante e con forti risvolti pratici.
	\end{itemize}

	\subsection{Fattori di rischio}

	\begin{itemize}
		\item Tecnologia reputata piuttosto complessa da padroneggiare adeguatamente;

		\item Possibili difficoltà nell'interazione con gli enti che andrebbero a certificare le competenze curricolari.
	\end{itemize}

	\subsection{Valutazione finale}

	Nonostante l'interesse suscitato dalle tecnologie proposte, l'elevata complessità del progetto e la dubbia efficacia del sistema di certificazione delle competenze curricolari da parte dell'ente erogatore ci spingono alla scelta di un altro capitolato.

	% CAPITOLATO 4

	\section{Capitolato C4}

	\begin{itemize}
		\item \textbf{Nome:} \textit{ECoRe}: Enterprise Content Recommendation;
		\item \textbf{Proponente:} SIAV SPA;
		\item \textbf{Commitente:} Prof.Tullio Vardanega e Prof. Riccardo Cardin.
	\end{itemize}

	\subsection{Descrizione}

	Realizzazione di un servizio proattivo in grado di suggerire all’utente, che accede a contenuti aziendali, altri contenuti di interesse che potrebbero essere utili nello svolgimento del proprio lavoro. Tale utilità sarà stabilita sulla base del comportamento dell’utente stesso.

	\subsection{Dominio applicativo}

	Il progetto si inserisce nell'ambito dell'\textit{enterprise search}, ovvero dell’applicazione delle conoscenze e dei metodi del reperimento dell’informazione, nel contesto della ricerca di informazioni all’interno di una organizzazione aziendale e in quello dei sistemi di raccomandazione, cioè strumenti software in grado di suggerire contenuti utili all’utente.

	\subsection{Dominio tecnologico}

	\begin{itemize}
		\item \textbf{Motore di ricerca:} \textit{Apache SolR o ElasticSearch};

		\item \textbf{Libreria di apprendimento automatico con sezione dedicata alla raccomandazione:} \textit{Apache Mahout};

		\item \textbf{Misure di sicurezza:} \textit{Identity and Access Management Keycloak};

		\item \textbf{Importazione di dati provenienti da fonti web:} \textit{Apache Nutch};

		\item \textbf{Documentazione:} \textit{Evernote}.
	\end{itemize}

	\subsection{Aspetti positivi}

	\begin{itemize}

		\item Esplorazione di diverse tecnologie interessanti.
	\end{itemize}
	\subsection{Fattori di rischio}

	\begin{itemize}
		\item Il numero di tecnologie richieste dalla realizzazione del progetto è molto elevato e la loro natura variegata;

		\item Lo use case del progetto è generico e non propriamente definito;

		\item Lo scopo del prodotto ha suscitato scarso interesse al gruppo.
	\end{itemize}

	\subsection{Valutazione finale}

	Lo scarso interesse suscitato e la generalità dello use case dello stesso ci spingono alla scelta di un altro capitolato.

	% CAPITOLATO 5

	\section{Capitolato C5}

	\begin{itemize}
		\item \textbf{Nome:} \textit{Ironworks};
		\item \textbf{Proponente:} Zucchetti Software srl;
		\item \textbf{Commitente:} Prof.Tullio Vardanega e Prof. Riccardo Cardin.
	\end{itemize}

	\subsection{Descrizione}

	Realizzazione di un software di costruzione di diagrammi di robustezza con la relativa generazione di codice Java per le entità persistenti.

	\subsection{Dominio applicativo}

	Il progetto si inserisce nell'ambito dell'automatizzazione della codifica allo scopo di aumentare la qualità e la velocità di produzione del codice.

	\subsection{Dominio tecnologico}

	\begin{itemize}
		\item \textbf{Disegnatore diagrammi UML:} \textit{ArgoUML, StarUML, Software Idea Designer} ecc.;

		\item \textbf{Codifica parte lato server:} \textit{Java} con server \textit{Tomcat} o \textit{Javascript} con server \textit{Node.Js};

		\item \textbf{Interfaccia grafica:} \textit{HTML5, CSS3, Javascript}.
	\end{itemize}

	\subsection{Aspetti positivi}

	\begin{itemize}
		\item Le tecnologie proposte per la realizzazione del progetto sono numerose ma abbordabili;

		\item Lo scopo del progetto è reputato molto utile e interessante.
	\end{itemize}

	\subsection{Fattori di rischio}

	\begin{itemize}
		\item Dubbi sulla capacità di garantire codice autogenerato di effettiva qualità.
	\end{itemize}

	\subsection{Valutazione finale}

	Il progetto proposto ha raccolto l'interesse della totalità dei membri del gruppo, tuttavia non tale da considerarla come prima scelta.

	% CAPITOLATO 6

	\section{Capitolato C6}

	\begin{itemize}
		\item \textbf{Nome:} \textit{Marvin}: dimostratore di Uniweb su Ethereum;
		\item \textbf{Proponente:} RedBabel;
		\item \textbf{Commitente:} Prof.Tullio Vardanega e Prof. Riccardo Cardin.
	\end{itemize}

	\subsection{Descrizione}

	L'obiettivo del progetto è la realizzazione di un sottoinsieme delle funzionalità di Uniweb come una Applicazione Decentralizzata che funzioni su di una \textit{Ethereum Virtual Machine}.

	\subsection{Dominio applicativo}

	Il progetto si inserisce nell'ambito di una sperimentazione della tecnologia BlockChain attraverso \textit{Ethereum}.

	\subsection{Dominio tecnologico}

	\begin{itemize}
		\item \textbf{Interfaccia grafica:} \textit{HTML5, CCS3, Javascript};

		\item \textbf{Interazione con \textit{Ethereum}:} \textit{Truffle11 development framework};

		\item \textbf{Misure di sicurezza:} \textit{Metamask}.
	\end{itemize}

	\subsection{Aspetti positivi}

	\begin{itemize}
		\item Le tecnologie proposte per la realizzazione del progetto sono interessanti, innovative e in forte sviluppo;
	\end{itemize}

	\subsection{Fattori di rischio}

	\begin{itemize}
		\item Il progetto stimola all'interno del gruppo un interesse relativamente basso a causa della mancata concretezza dello scopo del progetto, al di là dei fini dimostrativi in relazione alla tecnologia proposta;

		\item Le tecnologie necessarie alla realizzazione del progetto appaiono di un elevato grado di difficoltà di apprendimento.
	\end{itemize}

	\subsection{Valutazione finale}

	Nonostante l'interesse suscitato dalle tecnologie proposte, l'elevata complessità delle stesse e la mancanza di uno scopo concreto del prodotto ci spingono alla scelta di un altro capitolato.

	% CAPITOLATO 7

	\section{Capitolato C7}

	\begin{itemize}
		\item \textbf{Nome:} \textit{OpenAPM}: cruscotto di Application Performance Management;
		\item \textbf{Proponente:} Iks - Kirey Group;
		\item \textbf{Commitente:} Prof.Tullio Vardanega e Prof. Riccardo Cardin.
	\end{itemize}

	\subsection{Descrizione}

	Realizzare uno Strumento di APM basato su tecnologie \textit{Open Source}.

	\subsection{Dominio applicativo}

	Il progetto si inserisce nell'ambito degli \textit{APM(Application Performance Managment)}, strumenti per il monitoraggio e la gestione di performance ed availability delle applicazioni con l’obiettivo di individuare e diagnosticare in modo semplice problematiche complesse che impattano sul servizio erogato.

	\subsection{Dominio tecnologico}

	\begin{itemize}
		\item \textbf{Agent:} \textit{Java, PHP, Node.js, ElasticSearch};

		\item \textbf{Server:} \textit{Kibana, D3.js}.
	\end{itemize}

	\subsection{Aspetti positivi}

	\begin{itemize}
		\item Il progetto propone l'utilizzo di alcune tecnologie interessanti ed innovative;

		\item Realizzazione di un prodotto open source in un dominio applicativo caratterizzato da software proprietari.
	\end{itemize}

	\subsection{Fattori di rischio}

	\begin{itemize}
		\item Le numerose sfaccettature del progetto potrebbe portare al conseguimento di una qualità complessiva scadente.
	\end{itemize}

	\subsection{Valutazione finale}

	La complessità delle tecnologie da utilizzare e il timore di non realizzare un prodotto di qualità ci spingono alla scelta di un altro capitolato.

	% CAPITOLATO 8

	\section{Capitolato C8}

	\begin{itemize}
		\item \textbf{Nome:} \textit{TuTourSelf}: piattaforma di prenotazioni per artisti in tournee;
		\item \textbf{Proponente:} TuTourSelf S.r.l.;
		\item \textbf{Commitente:} Prof.Tullio Vardanega e Prof. Riccardo Cardin.
	\end{itemize}

	\subsection{Descrizione}

	Realizzare una piattaforma web (con focus particolare sull'esperienza utente) che permetta agli artisti indipendenti di tutto il mondo di organizzare in poco tempo il proprio tour comunicando direttamente con i locali disponibili.

	\subsection{Dominio applicativo}

	Il progetto si inserisce nell'ambito delle Single Page Applications (SPA), applicazioni web che possono essere usate o consultate su una singola pagina web con l'obiettivo di fornire una esperienza utente fluida e simile alle applicazioni desktop, con particolare attenzione all'esperienza utente.

	\subsection{Dominio tecnologico}

	\begin{itemize}
		\item \textbf{Front-end:} \textit{HTML5, CSS, Javascript, REACT};
		\item \textbf{Back-end:}  \textit{PHP, Ruby, Python, Java, MySQL}.
	\end{itemize}

	\subsection{Aspetti positivi}

	\begin{itemize}
		\item Lo scopo del progetto è reputato molto interessante;

		\item La difficoltà complessiva della realizzazione del progetto sembra essere relativamente bassa, date anche le conoscenze pregresse nell'ambito possedute dai membri del gruppo.
	\end{itemize}

	\subsection{Fattori di rischio}

	\begin{itemize}
		\item Il capitolato stimola all'interno del gruppo un interesse relativamente basso a causa delle sue dichiarate finalità di produrre una demo del prodotto vero e proprio, su cui raccogliere dati e fare analisi.
	\end{itemize}

	\subsection{Valutazione finale}

	Nonostante l'interesse suscitato dallo scopo del progetto e la sua evidente fattibilità, il suo effettivo utilizzo ci spinge alla scelta di un altro capitolato.

\end{document}
